% ==============================================================================
% SECCIÓN 10: VALOR AÑADIDO, BENEFICIOS Y BENEFICIARIOS
% ==============================================================================
\newpage
\section{Valor Añadido, Beneficios y Beneficiarios}

\subsection{Valor añadido de la propuesta}
Frente a las agendas visuales estáticas o diarios tradicionales analizados en el estudio previo (como \textit{In-TIC Agenda} o \textit{Día a Día}), \textbf{AUTOVIDA} incorpora capacidades orientadas a la autonomía activa, la gestión de imprevistos y la continuidad entre el aula y el entorno familiar:

\begin{center}
\small
\begin{tabularx}{\textwidth}{|>{\raggedright\arraybackslash}p{3.5cm}|>{\raggedright\arraybackslash}X|>{\raggedright\arraybackslash}p{3.2cm}|}
\hline
\rowcolor{tableHeader}
\textbf{Elemento diferenciador} & \textbf{Valor que aporta al centro y a los usuarios} & \textbf{Trazabilidad} \\
\hline
Gestión de imprevistos y actividades alternativas & Permite asociar una actividad de respaldo cuando la principal no puede completarse y reajustar la planificación si una tarea se prolonga, reduciendo el bloqueo o la frustración del alumnado ante cambios de rutina. & RF-13, RF-34 a RF-36 y Flujo principal \\
\hline
Iniciativa propia con control de prioridad & El alumnado puede crear sus propias tareas personales para fomentar su autogestión, pero el sistema garantiza que tengan prioridad inferior para no solapar ni desplazar la planificación docente. & RF-17 a RF-19 \\
\hline
Coordinación por franjas horarias (Centro--Hogar) & Delimita automáticamente los permisos de asignación según el horario: el personal gestiona la jornada escolar y los tutores/familias el horario extraescolar, unificando la agenda sin conflictos de edición. & RF-07, RF-10 y RF-20 \\
\hline
Continuidad sin conexión (\textit{Offline-first}) & Permite al alumnado consultar su agenda previamente sincronizada y marcar pasos completados sin depender de la cobertura de red en cada aula, sincronizando el avance al recuperar la conexión. & RNF (Funcionamiento sin conexión y Consistencia) \\
\hline
Seguimiento granular y alertas de cierre diario & Sustituye la revisión manual presencial por el registro paso a paso (obligatorio o guiado), temporizador pausable y notificaciones automáticas al final del día sobre tareas no completadas. & RF-22, RF-24 a RF-29 \\
\hline
Avisos ante cambios relevantes del alumno & Notifica a tutores y personal autorizados cuando se produzcan cambios en los datos, los apoyos o la planificación de un alumno que requieran su atención, facilitando una respuesta coordinada y evitando avisos innecesarios. & RF-41 \\
\hline
\end{tabularx}
\end{center}

\begin{cambiocontenido}{cambioBeneficios}
\subsection{Beneficios y beneficiarios}
Los beneficios esperados se distribuyen entre quienes utilizarán AUTOVIDA de forma directa y quienes pueden beneficiarse de una mejor coordinación educativa. Su alcance dependerá de los perfiles que el centro autorice, de la validación del MVP y de la incorporación de la herramienta a la dinámica habitual. La aplicación se plantea como apoyo a la intervención educativa y familiar, no como sustituto del acompañamiento personal ni del criterio profesional.

\begin{small}
\begin{longtable}{|>{\raggedright\arraybackslash}p{3.2cm}|>{\raggedright\arraybackslash}p{2.4cm}|>{\raggedright\arraybackslash}p{9.6cm}|}
\hline
\rowcolor{tableHeader}
\textbf{Beneficiario} & \textbf{Alcance previsto} & \textbf{Beneficio esperado y método de verificación} \\
\hline
\endfirsthead
\hline
\rowcolor{tableHeader}
\textbf{Beneficiario} & \textbf{Alcance previsto} & \textbf{Beneficio esperado y método de verificación} \\
\hline
\endhead
Alumnado de PTVAL (\textit{beneficiario directo y principal}) & Alumnado de las aulas PTVAL seleccionadas que participe en la validación. & \textbf{Beneficio:} Disponer de una agenda diaria más previsible y comprensible, consultar actividades con apoyos ajustados a cada perfil (texto, pictogramas, imágenes o recursos audiovisuales) y avanzar por instrucciones paso a paso. Las actividades alternativas y el registro del progreso pueden facilitar la gestión de cambios y favorecer una participación más autónoma, siempre con el acompañamiento que cada alumno necesite.\newline \textbf{Verificación:} Pruebas de aceptación y usabilidad acordadas con el centro: observar si los participantes localizan su actividad, comprenden los apoyos, marcan instrucciones, usan el temporizador y consultan una actividad alternativa; registrar dificultades y necesidades de apoyo (RF-13, RF-20 a RF-28; RNF-01, RNF-05 y RNF-06). \\
\hline
Profesorado (\textit{beneficiario directo}) & Docentes autorizados que planifican actividades y rutinas para el alumnado asociado. & \textbf{Beneficio:} Organizar las actividades escolares en un único espacio, reutilizar rutinas y consultar el avance registrado para preparar apoyos o ajustes posteriores. La asignación por horario y las alertas de cierre pueden facilitar el seguimiento, sin presuponer que reduzcan por sí solas el tiempo de preparación o supervisión.\newline \textbf{Verificación:} Completar escenarios de creación y asignación de actividades dentro del horario escolar, asignación de rutinas, consulta del progreso y recepción de avisos de tareas pendientes; recoger incidencias y valoración del profesorado (RF-10, RF-12, RF-22 a RF-29). \\
\hline
Personal de apoyo educativo (\textit{beneficiario directo}) & Profesionales autorizados que acompañan al alumnado durante la jornada y tienen alumnos asociados. & \textbf{Beneficio:} Consultar las actividades y apoyos previstos, registrar el progreso observado y coordinar la atención con el profesorado y los tutores autorizados. Esto puede ayudar a mantener una referencia común sobre la tarea en curso y los cambios planificados.\newline \textbf{Verificación:} Probar con cuentas autorizadas la consulta de alumnos asociados, el acceso a sus actividades, el registro del progreso y la recepción de notificaciones; comprobar también que no se accede a información de alumnos no asociados (RF-11, RF-22 a RF-29 y RNF-03). \\
\hline
Tutores legales y familias (\textit{beneficiarios directos y secundarios}) & Familias autorizadas vinculadas al alumnado participante. & \textbf{Beneficio:} Consultar y, dentro de los permisos y horarios definidos, organizar actividades del entorno extraescolar con apoyos coherentes con los utilizados en el centro. Las notificaciones sobre tareas pendientes o cambios relevantes pueden facilitar el seguimiento y la comunicación, sin convertir la aplicación en un canal de supervisión permanente.\newline \textbf{Verificación:} Comprobar la asignación y modificación de actividades fuera del horario escolar, la consulta de alumnos asociados y la recepción de avisos pertinentes; confirmar que usuarios no autorizados no acceden a los datos (RF-07 a RF-09, RF-11, RF-29 y RF-41; RNF-03). \\
\hline
Orientación y coordinación pedagógica (\textit{beneficiarios indirectos, si participan}) & Profesionales que colaboren en la definición de apoyos y en la validación acordada con el centro. & \textbf{Beneficio:} Contar con un recurso común para revisar si la presentación de actividades y los apoyos configurados resultan adecuados para las necesidades observadas, e identificar aspectos que convenga ajustar. AUTOVIDA no realiza diagnósticos ni sustituye la valoración profesional.\newline \textbf{Verificación:} Revisar con estos perfiles los ajustes de accesibilidad y los escenarios de uso, documentar observaciones y comprobar que los cambios acordados se reflejan en la configuración validada (RF-37, cuyo alcance está pospuesto, y RNF-01, RNF-05 y RNF-06). \\
\hline
Administración del sistema (\textit{beneficiario operativo}) & Personas designadas por el centro para gestionar las cuentas de alumnado, tutores y personal. & \textbf{Beneficio:} Disponer de operaciones centralizadas para crear, consultar, modificar y eliminar cuentas, facilitando la organización inicial de usuarios y el mantenimiento de la información básica del sistema.\newline \textbf{Verificación:} Ejecutar los casos de prueba de gestión de cuentas para cada tipo de usuario y comprobar que los cambios se reflejan correctamente y que se respetan los permisos definidos (RF-30 a RF-33 y RNF-03). \\
\hline
C.E.E. Fundación Purísima Concepción (\textit{beneficiario institucional}) & Comunidad educativa y responsables del programa PTVAL del centro. & \textbf{Beneficio:} Participar en la definición y validación de una solución ajustada a sus necesidades, disponer de documentación del producto y valorar su utilidad antes de decidir si se integra en sus procesos. Su sostenibilidad dependerá de los recursos de mantenimiento, soporte y protección de datos que se acuerden.\newline \textbf{Verificación:} Revisar y aceptar, según los criterios acordados, el MVP, la guía de uso y el informe de pruebas; registrar decisiones, limitaciones y propuestas de mejora al cierre de las iteraciones previstas (OBJ-04 y OBJ-05). \\
\hline
Futuras promociones de PTVAL (\textit{beneficiarios potenciales indirectos}) & Alumnado y profesionales que podrían utilizar la solución en cursos posteriores, si el centro decide mantenerla. & \textbf{Beneficio:} Aprovechar las lecciones aprendidas durante el piloto y, si se confirma la adecuación de la herramienta, contar con una base reutilizable que pueda adaptarse a nuevos perfiles y necesidades. Este beneficio no se considera alcanzado hasta que el centro valide la solución y acuerde su continuidad.\newline \textbf{Verificación:} Evaluar los resultados del piloto, documentar mejoras pendientes y comprobar si existe una decisión del centro sobre mantenimiento, adaptación y uso en cursos posteriores. \\
\hline
Equipo del proyecto y comunidad universitaria (\textit{beneficiarios académicos indirectos}) & Estudiantes y docentes que participan en el proyecto de Aprendizaje-Servicio. & \textbf{Beneficio:} Aplicar conocimientos de ingeniería del software, accesibilidad, gestión de requisitos y validación en colaboración con una entidad real, incorporando la retroalimentación del centro al trabajo académico.\newline \textbf{Verificación:} Revisar la trazabilidad entre necesidades, requisitos, prototipo y pruebas, así como las actas, entregables y reflexiones del proyecto. \\
\hline
\end{longtable}
\end{small}

El número exacto de cuentas y perfiles activos se acordará con la dirección del centro al inicio de la primera iteración. Para valorar los beneficios sin dar por logrados resultados que aún deben comprobarse, se establecerá una línea base y se compararán los resultados de las pruebas de validación. Entre los indicadores posibles se encuentran la proporción de actividades completadas con distintos niveles de apoyo, la identificación de actividades por parte del alumnado, el tiempo necesario para preparar una rutina y la valoración de familias y profesionales sobre la continuidad entre centro y hogar. La recogida de estos datos deberá acordarse con el centro, respetar la privacidad del alumnado y considerar el tamaño y las características de la muestra.
\end{cambiocontenido}
